\section{The complete Weil form in logarithmic coordinates}\label{sec:weil}
\subsection{Conventions and the classical criterion}
All Hilbert inner products are conjugate-linear in the first argument.
For $L>0$, let $I_L=(-L/2,L/2)$ and let $E_L$ denote extension by zero
from $I_L$ to $\R$. Unless another domain is specified, $f$ belongs to
$C_c^\infty(I_L;\C)$ and $F=E_Lf$. Our Fourier and translation
conventions are
\begin{equation}\label{eq:fourier}
 \widehat F(\tau)=\int_\R e^{-i\tau x}F(x)\dd x,
 \qquad U_dF(x)=F(x-d),\qquad
 \norm F^2=\frac1{2\pi}\int_\R|\widehat F(\tau)|^2\dd\tau.
\end{equation}
Write $\widetilde F(x)=\overline{F(-x)}$ and $g=F*\widetilde F$.
Then $g(0)=\norm F^2$, $g(-r)=\overline{g(r)}$, and
$g(-r)=\ip F{U_rF}$. In particular, the correlations vanish for $|r|\geq L$.

For $r>0$ set
\begin{equation}\label{eq:gamma-kernel}
 \gammak(r)=\frac{e^{-r/2}}{1-e^{-2r}}
 =\sum_{n\geq0}e^{-a_nr},\qquad a_n=2n+\tfrac12.
\end{equation}
One normalization of the additive Weil functional is
\begin{align}
 \mathcal W(g)={}&\int_\R 2\cosh(x/2)g(x)\dd x
 -\sum_{p,\,m\geq1}(\log p)p^{-m/2}
       \{g(m\log p)+g(-m\log p)\}\notag\\
 &-(\log(4\pi)+\gamma_E)g(0)\notag\\
 &-\int_0^\infty
       [g(r)+g(-r)-2e^{-r/2}g(0)]\gammak(r)\dd r.
 \label{eq:weil-functional}
\end{align}
Here $p$ runs over rational primes and $\gamma_E$ is Euler's constant.
The integrand at the origin has a finite limit, and the prime sum is
finite for compact support. This normalization can be found in
Suzuki's account of the Weil form \cite[Section 1.1]{Suzuki}; for its
geometric and trace interpretation see \cite{CCM}.

The Riemann hypothesis (RH) asserts that every nontrivial zero of the
meromorphic continuation of $\zeta(s)=\sum_{n\geq1}n^{-s}$ has real
part $1/2$. We use the following classical theorem as the number-theoretic
input; we do not reprove the explicit formula or the Weil criterion.
\begin{theorem}[Classical Weil positivity criterion]\label{thm:weil}
RH holds if and only if $\mathcal W(F*\widetilde F)\geq0$ for every
complex $F\in C_c^\infty(\R)$.
\end{theorem}
Every compact support is contained in some $I_L$. Consequently it is
sufficient, and necessary, to establish this positivity for every $L>0$
and every test function in $I_L$.

\subsection{Positive gamma energy and the fixed contact}
Define
\begin{align}
 K[F]&=\frac12\int_\R\int_\R
    |F(x)-F(y)|^2\gammak(|x-y|)\dd x\dd y\notag\\
 &=\int_0^\infty[2g(0)-g(r)-g(-r)]\gammak(r)\dd r.
 \label{eq:gamma-energy}
\end{align}
This is nonnegative without any arithmetic hypothesis. Expanding
\eqref{eq:gamma-kernel} and applying Plancherel gives
\begin{align}
 K[F]&=\frac1{2\pi}\int_\R b(\tau^2)|\widehat F(\tau)|^2\dd\tau,
 \label{eq:gamma-fourier}\\
 b(\tau^2)&=\sum_{n\geq0}\frac2{a_n}
                 \frac{\tau^2}{a_n^2+\tau^2}
 =\re\psi(\tfrac14+i\tau/2)-\psi(\tfrac14),
 \label{eq:gamma-multiplier}
\end{align}
where $\psi=\Gamma'/\Gamma$. Indeed,
$2\int_0^\infty(1-\cos(\tau r))e^{-ar}\dd r
=(2/a)\tau^2/(a^2+\tau^2)$; nonnegativity justifies interchange with
the sum. The second identity in \eqref{eq:gamma-multiplier} is the
digamma partial-fraction expansion \cite{DLMF}.

The archimedean part of \eqref{eq:weil-functional} equals
$K[F]+w_0\norm F^2$, where
\begin{align}
 w_0&= -\log(4\pi)-\gamma_E
       -2\int_0^\infty(1-e^{-r/2})\gammak(r)\dd r\notag\\
 &=\psi(\tfrac14)-\log\pi
 =-\gamma_E-\frac\pi2-3\log2-\log\pi<0.
 \label{eq:contact}
\end{align}
To verify the constant, the integral in the first line is
$\sum_{n\geq0}(a_n^{-1}-(2n+1)^{-1})
=\tfrac12[\psi(1/2)-\psi(1/4)]$.
The constant is therefore fixed by the arithmetic normalization.
It is not a parameter available to adjust a positive source.

\subsection{The signed poles and the localized target}
The first term of \eqref{eq:weil-functional} has kernel
$2\cosh((x-y)/2)$. Using the hyperbolic addition formula gives
\begin{equation}\label{eq:poles}
 P_L[f]=2\left|\int_{I_L}f(x)\cosh(x/2)\dd x\right|^2
       -2\left|\int_{I_L}f(x)\sinh(x/2)\dd x\right|^2.
\end{equation}
It is a rank-two Hermitian form, with one positive and one negative
direction for $L>0$. Its pointwise positive kernel does not make it a
positive quadratic form.

The exact target for this paper is
\begin{align}
 Q_L[f]:=\mathcal W(F*\widetilde F)
 ={}&K[F]+w_0\norm f^2+P_L[f]\notag\\
 &-2\sum_{m\log p<L}(\log p)p^{-m/2}
           \re\ip F{U_{m\log p}F}.
 \label{eq:target}
\end{align}
Its polarization is denoted by $Q_L(f,g)$. At equality $m\log p=L$
the compressed correlation is zero. The coefficient is the primitive
length $\log p$, not the repeated length $m\log p$.

On a fixed interval the prime and pole terms define bounded operators.
The gamma multiplier satisfies
\begin{equation}\label{eq:gamma-growth}
 b(\tau^2)=\log|\tau|-\log2-\psi(\tfrac14)+O(|\tau|^{-2})
 \quad (|\tau|\longrightarrow\infty).
\end{equation}
The natural energy domain is
\begin{equation}\label{eq:energy-domain}
 \mathcal D_{\gamma,L}=
 \left\{f\in L^2(I_L):
   \int_\R\log(2+|\tau|)|\widehat{E_Lf}(\tau)|^2\dd\tau<\infty\right\}.
\end{equation}
The positive gamma energy is closed on this domain; it is the restriction
of the closed nonnegative Fourier-multiplier form to the closed subspace
of functions supported in $\overline I_L$. The bounded other terms give
a closed semibounded form on the same domain. All matching statements
below are first required on smooth compactly supported inputs. The Weil
criterion itself needs only that test space.

Away from the diagonal the distribution kernel of $K$ is
$-\gammak(|x-y|)$, which is smooth whenever $x\ne y$. The prime
contribution instead has atoms at $x-y=\pm m\log p$. This distinction
will provide exact tests for candidate preparations.

Three small intervals already distinguish several issues:
\begin{center}
\begin{tabular}{lll}
\toprule
$L$ & Active prime powers & Test \\
\midrule
$1$ & $2$ & First contact, pole and translation identity \\
$5/4$ & $2,3$ & Joint source and mixed returns \\
$3/2$ & $2,3,4$ & First nontrivial repetition of $2$ \\
\bottomrule
\end{tabular}
\end{center}
A calculation on any one interval is preliminary. It does not establish
the all-support statement of Theorem~\ref{thm:weil}.

\subsection{The spectral form of the target}\label{subsec:spectral}
The target \eqref{eq:target} has a classical spectral description. We record
it because it governs what a realization has to produce, and because it
controls the size of the quantity being matched.

For $g=F*\widetilde F$ with $F\in C_c^\infty(\R)$, the function
\begin{equation}\label{eq:entire-transform}
 G(\tau)=\int_\R g(x)e^{-i\tau x}\dd x
 =\widehat F(\tau)\overline{\widehat F(\bar\tau)}
\end{equation}
is entire of exponential type and agrees with $|\widehat F(\tau)|^2$ for real
$\tau$. Write $\rho$ for the nontrivial zeros of $\zeta$ with multiplicity and
put $\gamma_\rho=-i(\rho-\tfrac12)$. In the normalization
\eqref{eq:weil-functional}, Weil's explicit formula states that
\begin{equation}\label{eq:explicit-formula}
 \mathcal W(g)=\sum_\rho G(\gamma_\rho),
\end{equation}
the sum taken in the symmetric order $|\im\rho|\leq T$, $T\to\infty$
\cite{Suzuki,CCM}. The first term of \eqref{eq:weil-functional} is
$G(i/2)+G(-i/2)$, the contribution of the two poles of $\zeta(s)$ at $s=0$
and $s=1$; the remaining terms are the archimedean and finite local factors.
We use \eqref{eq:explicit-formula} as classical input and do not reprove it.
Appendix~\ref{app:checks} records a direct check that the normalization
\eqref{eq:weil-functional} used here reproduces it.

\begin{corollary}[Spectral form of the localized target]\label{cor:spectral}
For $f\in C_c^\infty(I_L)$ and $F=E_Lf$,
\begin{equation}\label{eq:spectral-target}
 Q_L[f]=\sum_\rho\widehat F(\gamma_\rho)
        \overline{\widehat F(\overline{\gamma_\rho})}.
\end{equation}
Under RH every $\gamma_\rho$ is real and each summand is
$|\widehat F(\gamma_\rho)|^2\geq0$. A zero off the critical line pairs with
the zero $1-\bar\rho$ to contribute
$2\re[\widehat F(\gamma)\overline{\widehat F(\bar\gamma)}]$, which has no
fixed sign.
\end{corollary}

Three consequences organize what follows.

First, \eqref{eq:spectral-target} determines a realization that is
compatible in $L$. Suppose $\Phi$ is defined on $C_c^\infty(\R)$ and
satisfies $\norm{\Phi(f)}^2=Q_L[f]$ whenever $\supp f\subset I_L$. Then RH
holds, and $\Phi(f)\mapsto(\widehat F(\gamma_\rho))_\rho$ extends to a
unitary map from the closure of the range of $\Phi$ onto the closed span of
those vectors in $\ell^2$ of the zero multiset. The Hilbert space and the
source of a compatible realization are therefore not free data: up to that
unitary they are the zeros and the evaluation map.

Second, this fixes the shape of the answer without deciding the objective of
this manuscript. That some positive Hilbert space carries this spectrum is
equivalent to RH and cannot be arranged by fiat. The proposal here is more
specific: that an independently motivated family of positive pairings ---
protected line and defect algebras with physically defined twisted traces
--- contains such a source, so that construction of a named theory with
specified observables would supply it. Nothing in
\eqref{eq:spectral-target} decides whether that is so.

Third, \eqref{eq:spectral-target} is a warning about method. Its right side
is generally far smaller than the individual terms of \eqref{eq:target}. For
the smooth input used in Appendix~\ref{app:checks} on $I_4$, the gamma,
contact, pole and prime terms of \eqref{eq:target} are
$3.40$, $-8.06$, $9.11$ and $-4.46$, while their sum is $1.4\times10^{-7}$.
A candidate is therefore not usefully tested by matching those terms one at
a time. The cancellation between them is the content of the problem, and
Section~\ref{subsec:interference} turns that remark into an exact
obstruction.

\subsection{Parity and the hyperbolic pole pair}\label{subsec:parity}
The rank-two form \eqref{eq:poles} is usually read through its positive and
negative directions. A different basis makes its structure, and the
constraint it places on a realization, immediate.

\begin{lemma}[The pole form is a hyperbolic pair]\label{lem:polepair}
For $f\in C_c^\infty(I_L)$ and $F=E_Lf$ put
\begin{equation}\label{eq:polecoords}
 u[f]=\widehat F(i/2)=\int f(x)e^{x/2}\dd x,
 \qquad
 v[f]=\widehat F(-i/2)=\int f(x)e^{-x/2}\dd x,
\end{equation}
the evaluations at the two poles of $\zeta$, at $s=1$ and $s=0$. Then
$\ip{\cosh(x/2)}f=(u+v)/2$, $\ip{\sinh(x/2)}f=(u-v)/2$, and
\begin{equation}\label{eq:hyperbolic}
 P_L[f]=2\re\big(u[f]\,\overline{v[f]}\big).
\end{equation}
In the coordinates $(u,v)$ the form is therefore off-diagonal, with matrix
$\left(\begin{smallmatrix}0&1\\1&0\end{smallmatrix}\right)$: it is the
hyperbolic form of signature $(1,1)$, and its two null lines are $\{u=0\}$
and $\{v=0\}$.
\end{lemma}
\begin{proof}
$e^{\pm x/2}=\cosh(x/2)\pm\sinh(x/2)$ gives the two moment identities, and
$2|\tfrac{u+v}2|^2-2|\tfrac{u-v}2|^2=2\re(u\bar v)$. The matrix statement is
\eqref{eq:hyperbolic} read as a Hermitian form in $(u,v)$; the diagonal
entries vanish, so each coordinate axis is null and the signature is $(1,1)$.
\end{proof}

The negative direction of $P_L$ is thus not attached to either pole: it is the
antisymmetric combination, and the indefiniteness lies in the pairing between
the two poles rather than in one of them.

The same two combinations are the parity components. Let $(Rf)(x)=f(-x)$.
Every kernel in \eqref{eq:target} depends on $x-y$ through an even function,
so $Q_L(Rf,Rg)=Q_L(f,g)$, and the same holds for any form obtained from it by changing the
constant.
Both forms therefore decouple over the eigenspaces of $R$, and since
$\cosh(x/2)$ is even and $\sinh(x/2)$ odd,
\begin{equation}\label{eq:parity-split}
 P_L\big|_{\text{even}}=+2\big|\ip{\cosh(x/2)}f\big|^2,
 \qquad
 P_L\big|_{\text{odd}}=-2\big|\ip{\sinh(x/2)}f\big|^2 .
\end{equation}
The indefiniteness of the pole form is an odd-sector phenomenon and its
positive direction is entirely even.

\begin{proposition}[Every source carries the involution]\label{prop:involution}
Let $\Phi$ be defined on $C_c^\infty(I_L)$ with values in a Hilbert space and
$\norm{\Phi(f)}^2=Q_L[f]$. Then there is a unitary $J$ on the closure of the
span of the range of $\Phi$ with $J\Phi(f)=\Phi(Rf)$ and $J^2=\id$. The
positive direction of \eqref{eq:parity-split} lies in the $+1$ eigenspace of
$J$ and the negative direction in the $-1$ eigenspace. If moreover $\Phi$ is
compatible in $L$ in the sense of Corollary~\ref{cor:spectral}, then under the
unitary of that corollary $J$ is the involution $\rho\mapsto1-\rho$ of the
zero multiset.
\end{proposition}
\begin{proof}
Polarization gives $Q_L(Rf,Rg)=Q_L(f,g)$, so
$\ip{\Phi(Rf)}{\Phi(Rg)}=\ip{\Phi(f)}{\Phi(g)}$ and $\Phi(f)\mapsto\Phi(Rf)$
extends from the span of the range to an isometry, surjective because
$R^2=\id$; call it $J$. The eigenspace statement is \eqref{eq:parity-split}.
For the last statement, $\widehat{E_LRf}(\tau)=\widehat{E_Lf}(-\tau)$, and the
multiset $\{\gamma_\rho\}$ is invariant under $\gamma\mapsto-\gamma$ because
the functional equation pairs $\rho$ with $1-\rho$; so on the coordinates of
Corollary~\ref{cor:spectral} the map $J$ permutes $\gamma_\rho$ with
$\gamma_{1-\rho}=-\gamma_\rho$.
\end{proof}

The two poles $s=0$ and $s=1$ are exchanged by the same involution, being the
two points the functional equation adds to the zero set. Lemma~\ref{lem:polepair}
and Proposition~\ref{prop:involution} are used in
Section~\ref{subsec:leaves} to state what a realization of the pole term has
to contain.
